% !TeX root = ../main.tex

\chapter{Architettura del Sistema e Osservabilità}
\label{cap:architettura_osservabilita}

L'architettura del progetto è basata su \texttt{Kind} (\textit{Kubernetes IN Docker}), utilizzato per avviare cluster Kubernetes locali e leggeri all'interno di container Docker. Questa scelta consente di replicare molti comportamenti di un cluster reale (scheduler, networking, controller) in un ambiente contenuto e facilmente ricreabile, ideale per validazione e test prima di un eventuale rilascio in produzione.

Kubernetes è un orchestratore di container che gestisce deployment, scaling e resilienza delle applicazioni: organizza i container in \texttt{Pod}, mantiene lo stato desiderato tramite \texttt{Deployment} e \texttt{ReplicaSet}, espone i servizi con \texttt{Service} e fornisce funzionalità native di autoscaling e self-healing. Tutte queste funzionalità sono richiamate nel progetto tramite l'invocazione delle API di \texttt{kubectl}.

Docker è la piattaforma di containerizzazione che confeziona applicazioni e dipendenze in immagini portabili, garantendo isolamento, avvii rapidi e coerenza tra ambienti. Funge da base su cui costruire e distribuire le immagini eseguite da Kind (e quindi da Kubernetes).

\section{Design a Microservizi}
\label{sec:design_microservizi}

Il progetto adotta un design a microservizi che separa chiaramente responsabilità e domini di failure in tre blocchi indipendenti: il \textbf{Data Layer} (InfluxDB), dedicato alla persistenza di metriche e serie temporali; il \textbf{Message Broker} (MQTT/Mosquitto), che funge da canale asincrono di scambio eventi; e il \textbf{Logic Layer} (server centrale e servizi applicativi), che implementa l'elaborazione e l'esposizione delle funzionalità di business.

Questa separazione è anche fisica: ogni blocco viene eseguito in Pod distinti, con \texttt{Service} dedicati per la discovery via DNS interno (es. \texttt{influxdb-service} e \texttt{mosquitto-service}), in modo che il layer logico non dipenda da IP statici ma soltanto da endpoint di servizio. Il broker disaccoppia produttori e consumatori: il layer logico può scalare o riavviarsi senza spezzare l'ingest, mentre InfluxDB resta un componente \textit{stateful} con volume persistente (PVC) e una strategia di avvio che privilegia la consistenza dei dati.

Gli script Python non vengono eseguiti come processi manuali sui nodi: Kubernetes li avvia come processo principale del container tramite il \texttt{command} del Deployment (un container = un servizio, oppure più container nello stesso Pod quando serve co-locazione, ad esempio per condividere dati/volumi e ridurre la latenza). Lo scheduler decide su quale nodo eseguire ciascun Pod; il campo \texttt{replicas} permette di distribuire più istanze di uno stesso microservizio su nodi diversi per aumentare throughput e tolleranza ai guasti, mantenendo invariata l'interfaccia grazie ai \texttt{Service}.

\section{Simulazione del Mondo Fisico (Digital Twin)}
\label{sec:digital_twin}

I digital twin di droni e client sono implementati come generatori controllati di telemetria e domanda, utili a validare l'architettura end-to-end senza dipendere da dispositivi fisici.

Il twin del drone (in \texttt{drone\_simulator.py}) mantiene uno stato interno minimale ma realistico (posizione, batteria, usura, stato operativo) che evolve a tick: pubblica periodicamente la telemetria sul broker MQTT (topic \texttt{telemetry/drones}) e si mette in ascolto di comandi sul proprio canale dedicato \texttt{comandi/<DRONE\_ID>}, reagendo a eventi come l'assegnazione di una missione o il rientro alla base e simulando consumi e degrado fino alla condizione di \texttt{MAINTENANCE}.

In parallelo, il twin del client (in \texttt{client\_simulator.py}) genera carico di business pubblicando nuovi ordini su MQTT (topic \texttt{business/ordini/nuovi}) con intervalli casuali e payload strutturati (coordinate di pickup e drop, peso, priorità), così da stressare broker e layer applicativo in modo riproducibile.

Dal punto di vista infrastrutturale, entrambi vengono eseguiti su Kubernetes come Deployment separati, ciascuno nel proprio Pod (configurati nei manifest \texttt{configs/drone.yaml} e \texttt{configs/client.yaml}): il processo principale del container è lo script Python definito nel \texttt{command} del manifest, mentre la connettività verso l'infrastruttura è ottenuta via service discovery (\texttt{mosquitto-service}).

La moltiplicazione della flotta è ottenuta in modo nativo tramite lo scaling: aumentando il valore di \texttt{replicas} del Deployment \texttt{drone-simulator}, Kubernetes crea più Pod identici e, poiché \texttt{DRONE\_ID} viene derivato dal \texttt{metadata.name} del Pod, ogni replica corrisponde a un nuovo drone digitale con identità e canale comandi distinti. Lo scaling può inoltre essere attivato in modo automatico dall'agente: il server MCP espone il tool \texttt{scale\_drone\_deployment}, che invoca la patch della risorsa \texttt{deployment/scale} (equivalente concettuale di un \texttt{kubectl scale}), con i guardrail operativi descritti nel Capitolo~\ref{cap:safety_guardrails} per mantenere sotto controllo l'impatto delle azioni automatiche sul cluster.

\section{Osservabilità (Observability Strategy)}
\label{sec:osservabilita}

Un sistema distribuito deve permettere una chiara visibilità operativa per consentire il monitoraggio dello stato, la tracciabilità delle decisioni agentiche e la rapida mitigazione dei guasti. La strategia di osservabilità adottata si fonda su tre pilastri: la visualizzazione real-time tramite una dashboard web centralizzata, il tracciamento di log strutturati per garantire l'auditability delle azioni e la persistenza di metriche su serie temporali.

\subsection*{1. Dashboard Web Centralizzata}
La visibilità generale dello stato del sistema è affidata al server centrale (\texttt{central\_server.py}), che ospita un'interfaccia utente web integrata tramite template HTML asincroni e un set di endpoint REST nativi in Flask:

\begin{itemize}
    \item \texttt{/api/status}: fornisce uno snapshot numerico aggregato del sistema, con conteggio dei droni attivi, degli ordini pendenti e delle consegne completate, oltre alla configurazione dei nodi di rete (broker MQTT e istanza InfluxDB);
    \item \texttt{/api/orders}: espone la coda dei messaggi in stato \texttt{PENDING} memorizzati in RAM;
    \item \texttt{/api/drones}: restituisce la matrice completa dei droni agganciati al cluster, visualizzandone dinamicamente lo stato e le relative metriche;
    \item \texttt{/api/completed}: elenca lo storico logistico delle consegne andate a buon fine.
\end{itemize}

Lato client, la dashboard esegue un ciclo di polling asincrono tramite chiamate \texttt{fetch()} JavaScript ogni 5000\,ms (5 secondi). Questo intervallo garantisce un aggiornamento visivo costante senza saturare il server.

\subsection*{2. Log di Sistema e Audit Trail dei Comandi Agentici}
Il tracciamento dei log è strutturato su più livelli per garantire l'ispezionabilità (\textit{auditability}) richiesta dal dominio:

\begin{enumerate}
    \item \textbf{Registro di Audit Strutturato (\texttt{audit\_actions.jsonl}):} gestito nativamente dal modulo \texttt{drone\_mcp\_layer.py}, archivia in formato JSON Lines ogni singola operazione di scrittura o di rimediazione avviata dall'IA. Per ciascuna azione vengono tracciati il timestamp, il tipo di tool invocato e l'esatto payload utilizzato.
    \item \textbf{Registro delle Approvazioni Umane (\texttt{pending\_approvals.jsonl}):} traccia il ciclo di vita delle richieste degli agenti quando viene superato un confine di policy. Il file contiene la richiesta iniziale dell'agente, la giustificazione logica fornita dall'LLM e il successivo cambiamento di stato causato dall'operatore.
    \item \textbf{Log di Runtime e tracciamento dei costi:} i simulatori dei droni stampano a terminale log continui sul proprio stato. In parallelo, l'orchestratore e gli agenti loggano per ciascuna iterazione l'esatto consumo di token restituito dai modelli, metrica cruciale per la verifica dei tetti di spesa discussi nel Capitolo~\ref{cap:analisi_economica}.
\end{enumerate}

\subsection*{3. Metriche di Serie Temporali su InfluxDB}
L'ultimo pilastro dell'osservabilità riguarda l'archiviazione persistente dei dati operativi sotto forma di serie temporali. Quando il server centrale intercetta i messaggi sui topic MQTT, scrive in modalità sincrona record all'interno del bucket \texttt{iot\_data}:

\begin{itemize}
    \item \textbf{Measurement \texttt{drone\_telemetry}:} registra l'evoluzione fisica della flotta;
    \item \textbf{Measurement \texttt{business\_orders}:} monitora il flusso di ordini in ingresso.
\end{itemize}

Tramite il layer MCP (in particolare i tool \texttt{get\_drones\_telemetry} e \texttt{get\_pending\_orders}), gli LLM possono analizzare finestre temporali passate (\textit{time-windows} di 5 o 60 minuti) per estrarre trend di carico o pattern di deterioramento hardware, base informativa per le decisioni di scaling e manutenzione.
